\documentclass[11pt,a4paper]{article}
\usepackage{preamble}

\title{Lineer Cebir Notları}
\author{math-lab}
\date{\today}

\begin{document}
\maketitle

\section{Boyut ve rank}

\begin{definition}\label{def:rank}
  $T \colon V \to W$ lineer bir dönüşüm olsun. $T$'nin \emph{rank}'i,
  görüntü uzayının boyutudur: $\rank T = \dim \Img T$.
\end{definition}

Aşağıdaki sonuç lineer cebirin temel teoremlerinden biridir;
ayrıntılar için bkz.~\cite{axler2024}.

\begin{theorem}[Rank--nullity]\label{thm:rank-nullity}
  $V$ sonlu boyutlu bir vektör uzayı ve $T \colon V \to W$ lineer olsun. Bu durumda
  \begin{equation}\label{eq:rank-nullity}
    \dim V = \dim \Ker T + \rank T.
  \end{equation}
\end{theorem}

\begin{proof}
  $u_1, \dots, u_m$ vektörleri $\Ker T$'nin bir tabanı olsun. Bunları
  $v_1, \dots, v_n$ vektörleri ile $V$'nin bir tabanına genişletelim. Her $v \in V$ için
  \[
    v = \sum_{i=1}^{m} a_i u_i + \sum_{j=1}^{n} b_j v_j
    \quad\Longrightarrow\quad
    Tv = \sum_{j=1}^{n} b_j\, T v_j,
  \]
  dolayısıyla $\Img T = \Span(Tv_1, \dots, Tv_n)$ olur.

  Bu vektörlerin lineer bağımsız olduğunu gösterelim: $\sum_j c_j T v_j = 0$
  ise $T\bigl(\sum_j c_j v_j\bigr) = 0$, yani $\sum_j c_j v_j \in \Ker T$.
  Böylece bazı $d_i \in \R$ için $\sum_j c_j v_j = \sum_i d_i u_i$ yazılır;
  tabanın lineer bağımsızlığından tüm $c_j = 0$ bulunur.

  Sonuç olarak $\rank T = n$, $\dim \Ker T = m$ ve $\dim V = m + n$; bu da
  \cref{eq:rank-nullity}'i verir.
\end{proof}

\begin{corollary}\label{cor:injective}
  $\dim V = \dim W < \infty$ ise $T$ birebirdir $\iff$ örtendir.
\end{corollary}

\begin{proof}
  $T$ birebir $\iff \dim \Ker T = 0 \iff \rank T = \dim V = \dim W$
  (\cref{thm:rank-nullity}) $\iff T$ örten.
\end{proof}

\begin{remark}
  Bu teoremin Lean~4 ile formalleştirilmiş hâli için bkz.\\
  \nolinkurl{lean/MathLab/LinearAlgebra/RankNullity.lean}.
\end{remark}

\bibliographystyle{plain}
\bibliography{refs}

\end{document}
